% ============================================================
%  Schemes Reading Course — TU Berlin, Summer 2025
%  Following: Mumford, The Red Book of Varieties and Schemes
%
%  Contributing: compile with pdflatex or xelatex
% ============================================================

\documentclass[a4paper, 12pt]{article}

% ---- Packages -----------------------------------------------
\usepackage[utf8]{inputenc}
\usepackage[T1]{fontenc}
\usepackage{lmodern}
\usepackage{microtype}

\usepackage{amsmath, amsthm, amssymb}
\usepackage{mathtools}
\usepackage{enumitem}
\usepackage{hyperref}
\usepackage{geometry}
\geometry{margin=2.5cm}

% ---- Theorem environments -----------------------------------
\theoremstyle{plain}
\newtheorem{theorem}{Theorem}[section]
\newtheorem{proposition}[theorem]{Proposition}
\newtheorem{lemma}[theorem]{Lemma}
\newtheorem{corollary}[theorem]{Corollary}

\theoremstyle{definition}
\newtheorem{definition}[theorem]{Definition}
\newtheorem{example}[theorem]{Example}
\newtheorem{remark}[theorem]{Remark}

% ---- Common macros ------------------------------------------
\newcommand{\OO}{\mathcal{O}}          % structure sheaf
\newcommand{\Spec}{\operatorname{Spec}}
\newcommand{\Proj}{\operatorname{Proj}}
\newcommand{\kk}{k}                    % base field
\newcommand{\CC}{\mathbb{C}}
\newcommand{\ZZ}{\mathbb{Z}}
\newcommand{\QQ}{\mathbb{Q}}
\newcommand{\PP}{\mathbb{P}}
\newcommand{\AA}{\mathbb{A}}

% ---- Title --------------------------------------------------
\title{\textbf{Notes from the Schemes Reading Course}\\
\large TU Berlin, Summer Semester 2025\\
\normalsize Following Mumford, \emph{The Red Book of Varieties and Schemes}}

\author{Dominic Bunnett \and \textit{(add your name)}}

\date{Last updated: \today}

% =============================================================
\begin{document}
\maketitle
\tableofcontents
\newpage

% =============================================================
\section*{Introduction}
These are collaborative notes from the reading group on schemes at TU Berlin,
Summer Semester 2025. We are working through David Mumford's
\emph{The Red Book of Varieties and Schemes}~\cite{Mumford}, continuing from
where we left off last semester. Meetings take place on Wednesdays, 14:00--16:00,
in office EB~114.

Errors and improvements are welcome --- please submit a pull request or raise
an issue on the GitHub repository.

% =============================================================
\section{Varieties}

% Section stub — fill in as we go through Chapter I of the Red Book.

\subsection{Affine varieties}

\begin{definition}
  Let $\kk$ be an algebraically closed field. \emph{Affine $n$-space} over $\kk$
  is $\AA^n = \AA^n_\kk := \kk^n$.
\end{definition}

% Add further content here.

% =============================================================
\section{Schemes}

% Section stub — fill in as we go through Chapter II of the Red Book.

\subsection{Spec of a ring}

\begin{definition}
  Let $A$ be a commutative ring with unity. The \emph{prime spectrum} of $A$ is
  \[
    \Spec A := \{\, \mathfrak{p} \subset A \mid \mathfrak{p} \text{ is a prime ideal} \,\}.
  \]
  We equip $\Spec A$ with the \emph{Zariski topology}, in which the closed sets
  are of the form $V(\mathfrak{a}) := \{\mathfrak{p} \supset \mathfrak{a}\}$ for
  ideals $\mathfrak{a} \subset A$.
\end{definition}

% Add further content here.

% =============================================================
\bibliographystyle{alpha}
\begin{thebibliography}{9}

\bibitem{Mumford}
  D.~Mumford,
  \emph{The Red Book of Varieties and Schemes},
  Lecture Notes in Mathematics 1358,
  Springer, 2nd ed., 1999.

\end{thebibliography}

\end{document}
